\section{Билет 36. Уравнение состояния тела. Идеальный газ. Уравнение Менделеева - Клапейрона. Газовая постоянная.}

\begin{definition}
\textbf{Уравнением состояния} тела называется соотношение, связывающее между собой макроскопические параметры, характеризующие состояние данного вещества. Для простых газовых систем оно имеет вид:
\begin{equation}
    F(p, V, T) = 0, \label{eq:state_general_36}
\end{equation}
где $p$ --- давление, $V$ --- объём, $T$ --- абсолютная температура. Зная любые два параметра и уравнение состояния, можно однозначно определить третий.
\end{definition}

\begin{definition}
\textbf{Идеальным газом} называется теоретическая модель газа, в которой:
\begin{enumerate}
    \item Молекулы считаются материальными точками (их размеры пренебрежимо малы по сравнению с расстояниями между ними);
    \item Отсутствуют силы межмолекулярного взаимодействия, кроме моментов упругих столкновений;
    \item Столкновения молекул друг с другом и со стенками сосуда являются абсолютно упругими.
\end{enumerate}
Реальные газы при достаточно низких давлениях и высоких температурах с высокой точностью подчиняются законам идеального газа.
\end{definition}

\begin{theorem}[Уравнение Менделеева--Клапейрона]
Экспериментально установлено, что для фиксированной массы газа отношение произведения давления на объём к абсолютной температуре остаётся постоянным:
\begin{equation}
    \frac{pV}{T} = \text{const}. \label{eq:combined_gas_36}
\end{equation}
Константа пропорциональна массе газа. Если количество газа равно одному молю ($\nu = 1$), то объём называется \textbf{молярным} $V_m$, и величина становится универсальной для всех газов. Обозначим её через $R$ (универсальная газовая постоянная):
\begin{equation}
    p V_m = R T.
\end{equation}
Для произвольного количества вещества $\nu$ молей, занимающего объём $V = \nu V_m$, получаем \textbf{уравнение состояния идеального газа (уравнение Менделеева--Клапейрона)}:
\begin{equation}
    pV = \nu R T. \label{eq:men_clap_36}
\end{equation}
Учитывая, что $\nu = m/M$, где $m$ --- масса газа, а $M$ --- молярная масса, уравнение часто записывают в виде:
\begin{equation}
    pV = \frac{m}{M} R T.
\end{equation}
\end{theorem}

\begin{property}[Универсальная газовая постоянная]
\textbf{Газовая постоянная $R$} --- физическая константа, численно равная работе, совершаемой одним молем идеального газа при изобарном нагревании на один кельвин.
Её значение определяется из условий нормальных условий ($p_0 = 1.013 \cdot 10^5$ Па, $T_0 = 273.15$ К, $V_m = 22.414 \cdot 10^{-3}$ м³/моль):
\begin{equation}
    R = \frac{p_0 V_m}{T_0} \approx \frac{1.013 \cdot 10^5 \cdot 22.414 \cdot 10^{-3}}{273.15} \approx 8.314 \, \frac{\text{Дж}}{\text{моль} \cdot \text{К}}. \label{eq:R_value_36}
\end{equation}
В расчётах обычно используют $R \approx 8.31$ Дж/(моль·К).
\end{property}

\begin{remark}
Уравнение Менделеева--Клапейрона объединяет три частных газовых закона:
\begin{itemize}
    \item \textbf{Закон Бойля--Мариотта} ($T=\text{const}$): $pV = \text{const}$;
    \item \textbf{Закон Гей-Люссака} ($p=\text{const}$): $V/T = \text{const}$;
    \item \textbf{Закон Шарля} ($V=\text{const}$): $p/T = \text{const}$.
\end{itemize}
Связь универсальной газовой постоянной с постоянной Больцмана $k = 1.38 \cdot 10^{-23}$ Дж/К выражается через число Авогадро: $R = k N_A$. Это позволяет записать уравнение состояния через концентрацию молекул $n = N/V$:
\begin{equation}
    p = n k T,
\end{equation}
что является альтернативной формой уравнения идеального газа, связывающей макроскопическое давление с микроскопическими параметрами.
\end{remark}
